\documentclass[tikz,border=8pt]{standalone}
\usepackage[no-math]{fontspec}
\usepackage[brazil]{babel}
\usepackage{amsmath,amssymb}
\IfFileExists{C:/Windows/Fonts/Lato_Regular.ttf}{%
  \setmainfont{Lato}[
    Path = C:/Windows/Fonts/,
    UprightFont = Lato_Regular.ttf,
    BoldFont = Lato_Bold.ttf,
    ItalicFont = Lato_Italic.ttf,
    BoldItalicFont = Lato_BoldItalic.ttf
  ]
  \setsansfont{Lato}[
    Path = C:/Windows/Fonts/,
    UprightFont = Lato_Regular.ttf,
    BoldFont = Lato_Bold.ttf,
    ItalicFont = Lato_Italic.ttf,
    BoldItalicFont = Lato_BoldItalic.ttf
  ]
}{%
  \RequirePackage[default]{lato}
}
\renewcommand{\familydefault}{\sfdefault}

\usepackage{xcolor}
\usepackage{tikz}
\usepackage{pgfplots}
\pgfplotsset{compat=1.18}
\usetikzlibrary{arrows.meta,positioning,shapes.geometric,fit,backgrounds,calc,decorations.pathreplacing,patterns}

% Paleta institucional UEFS / CHF662
\definecolor{uefsNavy}{HTML}{1B2A4A}
\definecolor{uefsTeal}{HTML}{0E7C7B}
\definecolor{uefsForest}{HTML}{2D6A4F}
\definecolor{uefsGold}{HTML}{C4983B}
\definecolor{uefsGoldDark}{HTML}{8F6B1E}
\definecolor{uefsAlert}{HTML}{B02A37}
\definecolor{uefsOrange}{HTML}{D95F02}
\definecolor{uefsInk}{HTML}{22252A}
\definecolor{uefsMuted}{HTML}{5C6470}
\definecolor{uefsWater}{HTML}{2E86AB}
\definecolor{uefsSand}{HTML}{D9B382}
\definecolor{bgNavy}{HTML}{F2F5F9}
\definecolor{bgTeal}{HTML}{E9F5F4}
\definecolor{bgForest}{HTML}{EEF5F0}
\definecolor{bgGold}{HTML}{FCF7EB}
\definecolor{bgAlert}{HTML}{FDF2F2}
\definecolor{bgGray}{HTML}{F8FAFC}
\definecolor{lineGray}{HTML}{CFD8DC}

\tikzset{
  >=Stealth,
  aaCard/.style={draw=#1, fill=white, rounded corners=4pt, line width=1pt, inner sep=6pt, align=center, font=\small},
  aaBox/.style={draw=#1, rounded corners=4pt, line width=1.1pt, inner sep=6pt, align=center, font=\footnotesize\sffamily, text=uefsInk},
  aaTitle/.style={font=\normalsize\bfseries\sffamily, text=uefsNavy, align=center},
  aaSub/.style={font=\footnotesize\sffamily, text=uefsMuted, align=center},
  aaArrow/.style={->, >=Stealth, line width=1.4pt},
  aaDArrow/.style={<->, >=Stealth, line width=1.4pt}
}

\usepgfplotslibrary{fillbetween}

\begin{document}
\begin{tikzpicture}
\begin{axis}[
  width=14cm,
  height=7cm,
  scale only axis,
  title={Roteiro de Campo em Topossequência: Estações de Observação Integrada},
  title style={font=\normalsize\bfseries\sffamily, text=uefsNavy, yshift=6pt},
  axis lines=left,
  axis line style={lineGray, line width=1pt},
  tick style={uefsMuted},
  xmin=0, xmax=10.3,
  ymin=0, ymax=7.2,
  xtick={0, 2, 4, 6, 8, 10},
  ytick={0, 1, 2, 3, 4, 5, 6, 7},
  tick label style={font=\footnotesize\sffamily, text=uefsInk},
  xlabel={\textbf{Posição ao longo da topossequência (km / m relativos)}},
  ylabel={\textbf{Cota relativa (m)}},
  label style={font=\footnotesize\sffamily, text=uefsInk},
  legend style={
    at={(0.02,0.03)}, anchor=south west,
    font=\scriptsize\sffamily, fill=white, fill opacity=0.94,
    draw=lineGray, rounded corners=3pt, inner sep=4pt,
    cells={anchor=west}
  },
  legend image post style={scale=0.8},
  clip=true
]

% Superfície topográfica e substrato
\addplot[name path=base, draw=none, domain=0:10] {0};

\addplot[
  name path=bed,
  dashed, color=uefsNavy, line width=1pt,
  domain=0:10, samples=140, forget plot
] {
  (3.35 - 0.22*x - 1.0/(1 + exp(-(x - 4.2)*1.15)) + 0.65 - 0.32*exp(-((x - 9.55)^2)/0.22))
  - (0.68 - 0.30*exp(-((x - 4.2)^2)/2.0) + 0.22*exp(-((x - 8.9)^2)/1.4))
};

\addplot[
  name path=surf,
  color=uefsForest, line width=2.2pt,
  domain=0:10, samples=140, forget plot
] {3.35 - 0.22*x - 1.0/(1 + exp(-(x - 4.2)*1.15)) + 0.65 - 0.32*exp(-((x - 9.55)^2)/0.22)};

% Preenchimento substrato e manto de solo
\addplot[fill=lineGray!55, draw=none] fill between[of=base and bed];
\addlegendentry{Substrato rochoso / saprólito}

\addplot[fill=uefsSand!75, draw=none] fill between[of=bed and surf];
\addlegendentry{Manto de solo e colúvio}

% Re-traçar linhas para nitidez
\addplot[dashed, color=uefsNavy, line width=1pt, domain=0:10, samples=140, forget plot] {
  (3.35 - 0.22*x - 1.0/(1 + exp(-(x - 4.2)*1.15)) + 0.65 - 0.32*exp(-((x - 9.55)^2)/0.22))
  - (0.68 - 0.30*exp(-((x - 4.2)^2)/2.0) + 0.22*exp(-((x - 8.9)^2)/1.4))
};
\addplot[color=uefsForest, line width=2.3pt, domain=0:10, samples=140, forget plot]
  {3.35 - 0.22*x - 1.0/(1 + exp(-(x - 4.2)*1.15)) + 0.65 - 0.32*exp(-((x - 9.55)^2)/0.22)};

% Espelho d'água limpo no canal fluvial à direita (x = 9.08..9.98)
\addplot[name path=watertop, color=uefsWater, line width=1.1pt, domain=9.12:9.98, forget plot] {0.82};
\addplot[name path=waterbed, draw=none, domain=9.12:9.98, samples=40, forget plot]
  {3.35 - 0.22*x - 1.0/(1 + exp(-(x - 4.2)*1.15)) + 0.65 - 0.32*exp(-((x - 9.55)^2)/0.22)};
\addplot[fill=uefsWater!70, draw=none, forget plot] fill between[of=waterbed and watertop];
\draw[uefsWater, line width=1.1pt] (axis cs:9.12, 0.82) -- (axis cs:9.98, 0.82);
\node[font=\scriptsize\bfseries\sffamily, text=uefsWater] at (axis cs:9.55, 1.06) {Canal};

% 4 Estações acima do perfil topográfico com caixas escalonadas (zero sobreposição)
\node[aaBox=uefsForest, fill=bgForest, text width=2.75cm, inner sep=4pt, font=\scriptsize\sffamily]
  (p1) at (axis cs:1.30, 5.90)
  {\textbf{\color{uefsForest}P1 · Interflúvio / Topo}\\[2pt]
   Latossolo · infiltração · divisor};

\node[aaBox=uefsAlert, fill=bgAlert, text width=2.85cm, inner sep=4pt, font=\scriptsize\sffamily]
  (p2) at (axis cs:3.95, 5.20)
  {\textbf{\color{uefsAlert}P2 · Meia-Encosta}\\[2pt]
   Declividade · sulcos · transição A/Bt};

\node[aaBox=uefsGoldDark, fill=bgGold, text width=2.85cm, inner sep=4pt, font=\scriptsize\sffamily]
  (p3) at (axis cs:6.55, 4.15)
  {\textbf{\color{uefsGoldDark}P3 · Base / Colúvio}\\[2pt]
   Deposição · exfiltração · cobertura};

\node[aaBox=uefsNavy, fill=bgNavy, text width=2.85cm, inner sep=4pt, font=\scriptsize\sffamily]
  (p4) at (axis cs:8.85, 2.95)
  {\textbf{\color{uefsNavy}P4 · Canal / Zona Ripária}\\[2pt]
   APP ciliar Lei 12.651 · turbidez · margens};

% Linhas-guia verticais tracejadas + marcadores na superfície
\draw[dashed, color=uefsForest, line width=1pt] (p1.south) -- (axis cs:1.30, 3.72);
\filldraw[fill=uefsForest, draw=white, line width=0.8pt] (axis cs:1.30, 3.72) circle (3pt);

\draw[dashed, color=uefsAlert, line width=1pt] (p2.south) -- (axis cs:3.95, 2.70);
\filldraw[fill=uefsAlert, draw=white, line width=0.8pt] (axis cs:3.95, 2.70) circle (3pt);

\draw[dashed, color=uefsGoldDark, line width=1pt] (p3.south) -- (axis cs:6.55, 1.66);
\filldraw[fill=uefsGoldDark, draw=white, line width=0.8pt] (axis cs:6.55, 1.66) circle (3pt);

\draw[dashed, color=uefsNavy, line width=1pt] (p4.south) -- (axis cs:8.85, 1.03);
\filldraw[fill=uefsNavy, draw=white, line width=0.8pt] (axis cs:8.85, 1.03) circle (3pt);

\end{axis}

% Nota de rodapé abaixo do eixo
\node[font=\footnotesize\itshape\sffamily, text=uefsMuted, anchor=north]
  at (7.0, -1.05)
  {Registrar coordenadas GNSS, foto com escala, clinometria, solo, feições erosivas e APP.};

\end{tikzpicture}
\end{document}
